% Reference entries, J (AN-11).  Law 1: os/iword/core/edit.c iw_ed_just
% ([Just:] D0/06 at the cursor, a [DSRt] first when the cursor is not at the
% left margin; 0 Left, 1 Full (the default), 2 Center, 3 Right),
% spec/iword/wp5_format.h "margins and justification", face_a.c M_FMTJUST
% (IW_FMT_JUST_PROMPT).  Shots: ../wordperfect51-decomp mint MINT22, MINT24.
% GATES: test-pubs-iword test-iword-codes test-iword-wp5-roundtrip test-iword-face test-flair-iword-format
\Gates{test-pubs-iword test-iword-codes test-iword-wp5-roundtrip test-iword-face test-flair-iword-format}
\begin{Feature}{Justification}
Justification sets where lines sit between the margins: against the left
margin, centered, against the right margin, or full.  It applies from the
place you set it to the end of the document, or to the next Justification
code.

The four justifications are:
\begin{description}
\item[Left] every line starts at the left margin; the right edge is ragged.
\item[Center] every line is centered between the margins: use it for
headings and title lines.
\item[Right] every line ends at the right margin: use it for a date or a
return address.
\item[Full] the justification of a new document.  On the editing screen,
lines with full justification start at the left margin, as left-justified
lines do.
\end{description}

To set justification,
\begin{WPSteps}
\item Move the cursor to where the justification is to begin.
\item Press \FK{Format}{Shift-F8}, then select \Ch{Line}{1}.
\item Select \Ch{Justification}{3}.  The bottom line shows
\StatusLine{Justification: 1 Left; 2 Center; 3 Right; 4 Full: 0}
\item Select \Ch{Left}{1}, \Ch{Center}{2}, \Ch{Right}{3} or \Ch{Full}{4}.
The Format: Line screen shows the new justification.
\item Press \FK{Exit}{F7} to return to the document.
\end{WPSteps}
\Documents{wpmenu}{Shift-F8,1,3=Justification}
\Documents{wpmenu}{Shift-F8,1,3,1=Left}
\Documents{wpmenu}{Shift-F8,1,3,2=Center}
\Documents{wpmenu}{Shift-F8,1,3,3=Right}
\Documents{wpmenu}{Shift-F8,1,3,4=Full}

\BeginReplay{ref-just}
For example, to type a centered heading and a line against the right
margin, press \Press{Format}{Shift-F8}, select \Choose{Line}{1},
\Choose{Justification}{3} and \Choose{Center}{2}, press \Press{Exit}{F7},
type \TypeRun{Quarterly Report} and press \PressKey{Enter}.  Then press
\KC{Shift-F8}{1,3,3} and \K{F7}\Keys{Shift-F8}\Typing{133}\Keys{F7}, type \TypeRun{March 2, 1992} and press
\Press{Reveal Codes}{Alt-F3}:
\WordScreen[rows=14]{ref-just}{A centered line and a right-justified line, with their codes}
\EndReplay

\Notes
\Topic{Code}
Justification puts a \Code{[Just:Left]}, \Code{[Just:Center]},
\Code{[Just:Right]} or \Code{[Just:Full]} code in the document.  On the
editing screen InitechWord asks before deleting one (see \InRef{Codes}).
Delete it to return to the justification in effect before it.

\Topic{Justification in the Middle of a Line}
Justification takes effect at the start of a line.  If the cursor is not at
the left margin when you set it, InitechWord first ends the line at the
cursor with a deletable soft return, \Code{[DSRt]}, so that the text after
the cursor begins a new line (see \InRef{Return, Deletable Soft}).

\BeginReplay{ref-just-mid}
For example, type \TypeRun{Mid}, then press \PressKC{Shift-F8}{1,3,2}{132}
and \Press{Exit}{F7} with the cursor after it, type \TypeRun{after} and press
\Press{Reveal Codes}{Alt-F3}:
\WordScreen[rows=14]{ref-just-mid}{Justification set in the middle of a line}
\EndReplay

\Topic{Status Line}
Pos on the status line gives the cursor's place on the paper, so on a
centered or right-justified line it shows where the text really is.

\SeeAlso{Format; Margins, Left and Right; Return, Deletable Soft}
\end{Feature}
